%% Hand-authored circuitikz figure -- NOT generated by maestro2tex.
%% Abstracted symbol (sch_cell_symbol.tex), two leads: the single interface is a
%% two-terminal element. Internals deliberately not shown; values live in
%% note_models.tex.
\providecommand{\MaestroRoot}{include/analog/}
\input{\MaestroRoot sch_cell_symbol.tex}
\begin{figure}[htbp]
  \centering
\begin{circuitikz}[
    every node/.append style={font=\footnotesize},
    nlab/.style={font=\scriptsize, inner sep=1.5pt},
  ]
\CellSymTwo{0,0}{1.0}{\texttt{a}}{\texttt{b}}
\end{circuitikz}
  \caption[{Single-interface model.}]{Single electrode--electrolyte interface, presented as a two-terminal element and shown as the same abstracted symbol. Its elements and their values are given in Table~\ref{tab:analog-models}; they describe a much smaller double layer and a much larger charge-transfer resistance than the three-electrode cell. None of the benches reported in this chapter uses it.}
  \label{fig:randles-electrode}
\end{figure}
